% !TEX root = ../metatime_monograph.tex
\chapter{Anomaly Cancellation Conditions}
\label{ch:anomaly_cancellation}

\section{The Five Conditions}

The Standard Model requires cancellation of five gauge anomalies.
With Weyl fermion multiplicities included, the five conditions are:

\[
\begin{aligned}
\text{[SU(3)]}^2\text{U(1)}_Y &: 
2N_c Y_Q - N_c Y_{uR} - N_c Y_{dR} = 0, \\[4pt]
\text{[SU(2)]}^2\text{U(1)}_Y &: 
3(2Y_Q) + 2Y_L = 0, \\[4pt]
\text{[U(1)}_Y]^3 &: 
2N_c Y_Q^3 + 2Y_L^3 - \left(N_c Y_{uR}^3 + N_c Y_{dR}^3 + Y_{eR}^3\right) = 0, \\[4pt]
\text{[Grav]}^2\text{U(1)}_Y &: 
2N_c Y_Q + 2Y_L + N_c Y_{uR} + N_c Y_{dR} + Y_{eR} = 0, \\[4pt]
\text{[SU(2)]}^2\text{U(1)}_B &: 
3Y_Q - Y_{uR} - Y_{dR} = 0.
\end{aligned}
\]

The factors account for:
\begin{itemize}
\item \(N_c = 3\) colours for quarks.
\item Factor 2 for SU(2) doublets (left-handed fields).
\item Singlet counting for right-handed fields.
\end{itemize}

\section{Verification with Metatime Values}

\subsection*{Condition 1: \([\text{SU(3)}]^2\text{U(1)}_Y\)}

\[
2N_c Y_Q - N_c Y_{uR} - N_c Y_{dR}
= 6\!\left(\frac16\right) - 3\!\left(\frac23\right) - 3\!\left(-\frac13\right)
= 1 - 2 + 1 = 0.
\]

\subsection*{Condition 2: \([\text{SU(2)}]^2\text{U(1)}_Y\)}

\[
6Y_Q + 2Y_L
= 6\!\left(\frac16\right) + 2\!\left(-\frac12\right)
= 1 - 1 = 0.
\]

\subsection*{Condition 3: \([\text{U(1)}_Y]^3\)}

\[
\begin{aligned}
& 2N_c Y_Q^3 + 2Y_L^3 - \left(N_c Y_{uR}^3 + N_c Y_{dR}^3 + Y_{eR}^3\right) \\
&= 6\!\left(\frac{1}{216}\right) + 2\!\left(-\frac{1}{8}\right)
   - \left[3\!\left(\frac{8}{27}\right) + 3\!\left(-\frac{1}{27}\right) + (-1)\right] \\[4pt]
&= \frac{6}{216} - \frac{2}{8} - \left[\frac{24}{27} - \frac{3}{27} - 1\right] \\[4pt]
&= \frac{1}{36} - \frac{1}{4} - \left[\frac{21}{27} - 1\right] \\[4pt]
&= \frac{1}{36} - \frac{9}{36} - \left[\frac{7}{9} - \frac{9}{9}\right] \\[4pt]
&= -\frac{8}{36} - \left[-\frac{2}{9}\right] \\[4pt]
&= -\frac{2}{9} + \frac{2}{9} = 0.
\end{aligned}
\]

\subsection*{Condition 4: \([\text{Grav}]^2\text{U(1)}_Y\)}

\[
\begin{aligned}
& 2N_c Y_Q + 2Y_L + N_c Y_{uR} + N_c Y_{dR} + Y_{eR} \\
&= 6\!\left(\frac16\right) + 2\!\left(-\frac12\right)
   + 3\!\left(\frac23\right) + 3\!\left(-\frac13\right) + (-1) \\[4pt]
&= 1 - 1 + 2 - 1 - 1 = 0.
\end{aligned}
\]

\subsection*{Condition 5: \([\text{SU(2)}]^2\text{U(1)}_B\)}

\[
3Y_Q - Y_{uR} - Y_{dR}
= 3\!\left(\frac16\right) - \frac23 - \left(-\frac13\right)
= \frac12 - \frac23 + \frac13
= \frac{3 - 4 + 2}{6} = \frac16 \neq 0.
\]

\section{The Baryon Number Anomaly}

Condition 5 (baryon number anomaly \([\text{SU(2)}]^2\text{U(1)}_B\)) is
not cancelled in the SM because baryon number is anomalous.  The
non-zero value \(1/6\) is the standard baryon number anomaly coefficient
that gives rise to the electroweak sphaleron process.

The sphaleron rate is:

\[
\Gamma_{\text{sph}} \propto \exp\left(-\frac{4\pi}{\alpha_W}\right),
\]

which sets the scale for baryon number violation in the early universe.

\section{Summary}

All four gauge anomalies \(A_1\) through \(A_4\) cancel identically with
the Metatime-derived hypercharges.  The fifth condition (\(A_5\),
baryon number) does not cancel, consistent with the Standard Model
prediction of electroweak baryogenesis.
